\documentclass{article} % For LaTeX2e
\usepackage[submission]{colm2026_conference}

\usepackage{microtype}
\usepackage{hyperref}
\usepackage{url}
\usepackage{booktabs}
\usepackage{amsmath}
\usepackage{lineno}

\definecolor{darkblue}{rgb}{0, 0, 0.5}
\hypersetup{colorlinks=true, citecolor=darkblue, linkcolor=darkblue, urlcolor=darkblue}

\title{Honest Bits-per-Byte: Leak-Free Causal Patch Boundaries\\ for Hierarchical Byte-Level Language Models}

% Anonymized for submission via the [submission] option.
\author{Anonymous Authors}

\newcommand{\rg}{\textsc{root\_greedy}}
\newcommand{\bt}{\textsc{bt}}

\begin{document}
\maketitle

\begin{abstract}
Hierarchical byte-level language models (e.g.\ AU-Net, the Byte Latent Transformer) replace a
fixed token vocabulary with a learned hierarchy: raw bytes are \emph{pooled} into patches at a set
of patch boundaries, a trunk transformer runs on the pooled sequence, and a decoder upsamples back
to bytes. When those boundaries are derived from a reference BPE vocabulary, the standard
maximal-munch (backtracking longest-match) segmentation \emph{leaks future bytes} into the
teacher-forced training signal: deciding where patch $k$ ends consumes bytes belonging to patch
$k{+}1$. We show this leak inflates training bits-per-byte (BPB) by $\approx0.07$ at the 1.3B scale
--- a gain that is \emph{not realizable at generation}, where the decoder must produce its own
boundaries causally. We introduce \rg{}, a boundary scheme that is leak-free on both the
\emph{mode} axis (a greedy, prefix-only, $O(n)$ trie walk that never backtracks) and the
\emph{placement} axis (boundaries placed at the next patch's start, so predicting a byte never reads
the boundary it triggers). \rg{} makes training boundaries \emph{byte-identical} to the streaming
decoder's, giving an honest BPB of $0.860$ (vs.\ $0.787$ for the leaky variant), while retaining the
full byte-level robustness profile: it is competitive-to-best on reasoning, leads a subword Llama
baseline on character manipulation and typo robustness, and is \emph{exactly} invariant to prompt
cut-point shifts (the ``prompt boundary problem'').
\end{abstract}

\section{The leak in hierarchical byte models}
\vspace{-2pt}
Token-free byte models~\citep{xue2022byt5, slagle2024spacebyte} avoid a fixed vocabulary but pay an
$O(n)$ sequence-length cost; hierarchical byte
models~\citep{nawrot2023hourglass, yu2023megabyte, pagnoni2024blt, pagnoni2025aunet} recover
efficiency by pooling bytes into patches. Such a model is defined by a single segmentation function
$f:\text{bytes}\!\to\!\text{boundary positions}$. The tokenizer itself is trivial (raw bytes; vocab
$258 = 256 + \text{BOS} + \text{EOS}$); ``BPEByte'' means only that the patch boundaries are placed
where a reference BPE vocabulary~\citep{grattafiori2024llama3} would cut the byte stream. The boundaries determine (a) how bytes are pooled for the trunk and (b) via the
decoder's gather index \texttt{repeat\_idx}, \emph{which} patch state each byte may read when it is
predicted. If the boundary for byte $e$ is computed from bytes at or after $e$, teacher forcing hands
the model future information it cannot have at generation --- a \textbf{train/generation leak}.

There are two independent leak axes. (i)~\textbf{Mode.} Maximal-munch (\bt{}: backtracking
longest-match) walks the trie, remembers the longest \emph{valid} token, and at a dead-end commits it
and re-feeds the overshoot. The end of patch $k$ is thus a function of bytes \emph{after} patch $k$.
(ii)~\textbf{Placement.} With \emph{before\_root} placement the boundary sits at a token's last byte
($e{-}1$); since byte $e$ gathers \texttt{repeat\_idx}$[e{-}1]$, the boundary that $e$ itself triggers
is visible when predicting $e$.

\section{\rg{}: leak-free on both axes}
\vspace{-2pt}
\rg{} closes both leaks simultaneously (Table~\ref{tab:axes}).

\noindent\textbf{Greedy mode.} Walk the trie and commit \emph{all} accumulated bytes the instant the
walk hits a dead-end --- no longest-valid-token memory, no re-feed. The cut depends only on the bytes
walked so far, is $O(n)$ in a single left-to-right pass, and \emph{never rolls back} --- exactly the
forward-only commitment a streaming decoder makes. (The committed patch need not be a valid vocab
token; patches are pooling spans, not units to be decoded.)

\noindent\textbf{Root placement.} Put the boundary at the \emph{next} patch's first byte (index $e$,
offset $0$) rather than at $e{-}1$. Then byte $e$'s gather \texttt{repeat\_idx}$[e{-}1]$ structurally
\emph{excludes} the boundary at $e$: predicting $e$ cannot see the decision $e$ triggers.

Because greedy never backtracks, its streaming frontier is monotone and equals the batched result, so
the \textbf{train/generation boundary gap is exactly $0$} (verified byte-identical against the
streaming parser). \rg{} is the only mode$\times$placement combination that is leak-free \emph{and}
causal on both axes.

\begin{table}[t]
\centering
\small
\begin{tabular}{lccccc}
\toprule
scheme & mode & placement & lookahead to cut? & boundary visible? & net \\
\midrule
offline \bt{} & longest-match & before\_root & yes & yes & leaky ($\times2$) \\
online \bt{}  & longest-match$+$backtrack & before\_root & yes & yes & leaky \\
\textbf{\rg{}} & \textbf{greedy} & \textbf{root} & \textbf{no} & \textbf{no} & \textbf{leak-free $+$ causal} \\
\bottomrule
\end{tabular}
\caption{The two leak axes. \rg{} is the only combination leak-free on both.}
\label{tab:axes}
\vspace{-10pt}
\end{table}

\section{Results}
\vspace{-2pt}
We train 1.3B-parameter models on DCLM~\citep{li2024dclm} for 180k steps ($\approx283$\,GB of bytes),
and compare against a subword Llama-1.8B baseline (paper recipe, iso-byte at 283\,GB) and AU-Net2
(pure-byte whitespace patches). BPB is tokenizer-agnostic (EMA, $\alpha{=}0.01$).

\noindent\textbf{The leak is worth $\approx0.07$ BPB and is not real
(Table~\ref{tab:bpb}).} The leaky \bt{} variants ($0.78$--$0.79$) sit far below everything, but the
gap is purely boundary lookahead the model cannot use when decoding. \rg{} (0.860) is the
\emph{honest} number; it $\approx$ AU-Net2 (0.866), i.e.\ two different leak-free byte segmentations
land in the same place, and both sit just above the subword baseline (0.839). Crucially, the
$0.073$ gap between online-\bt{} and \rg{} comes from \emph{nothing but the leak} --- same bytes,
same architecture.

\begin{table}[t]
\centering
\small
\begin{minipage}[t]{0.46\textwidth}
\centering
\begin{tabular}{lc}
\toprule
model (1.3B, 283\,GB) & train BPB \\
\midrule
BPEByte offline-\bt{} (leaky) & 0.780 \\
BPEByte online-\bt{} (leaky)  & 0.787 \\
Llama 1.8B (subword)          & 0.839 \\
\textbf{BPEByte \rg{} (leak-free)} & \textbf{0.860} \\
AU-Net2 (pure byte)           & 0.866 \\
\bottomrule
\end{tabular}
\caption{Honest BPB. The leaky variants peek; \rg{} does not.}
\label{tab:bpb}
\end{minipage}\hfill
\begin{minipage}[t]{0.50\textwidth}
\centering
\begin{tabular}{lccc}
\toprule
metric (0-shot) & Llama & AU-Net2 & \rg{} \\
\midrule
HellaSwag $acc_n$  & 62.2 & 62.6 & \textbf{63.0} \\
ARC-Easy $acc_n$   & 65.4 & 65.7 & \textbf{65.9} \\
ARC-Chal.\ $acc_n$ & 35.3 & 36.5 & \textbf{36.6} \\
PIQA acc           & 74.9 & 73.0 & 73.7 \\
CUTE (14 tasks)    & 18.4 & \textbf{23.9} & 20.6 \\
\bottomrule
\end{tabular}
\caption{Downstream: \rg{} is competitive-to-best on reasoning and beats subword on CUTE.}
\label{tab:down}
\end{minipage}
\vspace{-12pt}
\end{table}

\noindent\textbf{The higher honest BPB does \emph{not} cost downstream quality
(Table~\ref{tab:down}).} Despite training BPB above the leaky variants, \rg{} is competitive-to-best
on reasoning (best of the three on HellaSwag, ARC-Easy, ARC-Challenge $acc_{norm}$) and keeps the
character-level strengths that motivate byte models: CUTE~\citep{edman2024cute} 20.6 vs.\ Llama 18.4, the smallest typo
degradation under noisy evaluation (BoolQ $+0.8$ pts under typos, vs.\ Llama $-4.0$), and the best
HellaSwag-Noise average (42.2 vs.\ 37.6).

\noindent\textbf{Exact cut-point invariance.} Byte models consume the same raw bytes regardless of
where a prompt is cut, so they are immune to the ``prompt boundary problem'' (the familiar
``never end your prompt with a trailing space''). On a cut-point BPB probe, Llama-1.8B shows
$\Delta\text{BPB}_{\text{en}}=+1.18$ (and flips MCQ answers, $\Delta\text{Acc}=-8.9\%$ when a trailing
space is committed), whereas \rg{} is invariant by construction: $\Delta\text{BPB}=0.000$ and MCQ
$\Delta\text{Acc}=+0.16\%$.

\section{Takeaway}
\vspace{-2pt}
When a hierarchical byte model's patch boundaries are derived from BPE, the segmentation function
enters the loss, and naive maximal-munch silently leaks the future. The resulting BPB improvement is
a \emph{diagnostic upper bound} --- how much a model \emph{could} gain from boundary lookahead --- not
a deployable result. \rg{} reports the honest number by forcing training boundaries to equal the
ones the model's own decoder produces, at $O(n)$ cost and with zero train/generation gap, while
preserving every byte-level robustness advantage. We recommend it as the default for any honest
bits-per-byte claim or any deployment where the model autoregressively produces its own boundaries,
and \bt{} only as a leakage diagnostic.

\bibliography{references}
\bibliographystyle{colm2026_conference}

\end{document}
